\subsection{The direction selected by the dodecaphase register}
\label{exc:k-direccion}

The record \(K\) has a geometric role that complements its rational
evaluation and its support readout. Once the representation of the twelve
positions has been constructed, its components determine a direction in the
space of zero-mean variations. The reduction \(12\to11\) removes the
uniform mode; the subsequent reduction \(11\to10\) removes the direction
selected by the record itself. The two operations therefore have different
antecedents.

To specify the chart used completely, consider the group \(A_5\) of even
permutations of five objects and the subgroup
\(C_5=\langle(1\,2\,3\,4\,5)\rangle\). The twelve left cosets
\(r_pC_5\) are ordered by the following representatives, written as image
lists:
\[
\begin{array}{c|c@{\qquad}c|c}
p&r_p&p&r_p\\ \hline
1&(4,3,5,2,1)&7&(5,4,3,2,1)\\
2&(4,2,1,3,5)&8&(1,3,4,2,5)\\
3&(1,2,4,5,3)&9&(4,2,3,5,1)\\
4&(2,5,3,4,1)&10&(3,2,5,4,1)\\
5&(4,3,1,5,2)&11&(4,5,2,3,1)\\
6&(5,1,2,3,4)&12&(5,3,2,4,1).
\end{array}
\]
The action by left multiplication defines an orthogonal representation
\(\rho:A_5\to O(\RR^{12})\): \(\rho(a)e_p=e_q\) whenever
\(ar_pC_5=r_qC_5\). This representation chart acts on the record already
obtained; the group characters are not used to choose its twelve values.

Let \(\chi_3\) denote the values of the three-dimensional character
\[
\begin{array}{c|ccccc}
\text{class}&1&2^2&3&5_a&5_b\\ \hline
\chi_3&3&-1&0&(1+\sqrt5)/2&(1-\sqrt5)/2.
\end{array}
\]
The class \(5_a\) contains the displayed generator of \(C_5\), and
\(5_b\) contains its square. This convention distinguishes the two conjugate
three-dimensional sectors of the representation. Define
\begin{equation}
 P_3=\frac3{60}\sum_{a\in A_5}\chi_3(a^{-1})\rho(a),
 \qquad P_{11}=I_{12}-\frac1{12}J_{12}.
 \label{eq:k-proyectores-incidencia}
\end{equation}
Both matrices are determined by finite sums in \(\QQ(\sqrt5)\). The
appearance of this quadratic field belongs to the realization of the
representation, subsequent to the regional generation of \(\varphi\).

\begin{lemma}[Projection of the record onto the three-dimensional sector]
\label{lem:k-sector-no-nulo}
One has
\[
 P_3=P_3^{\mathsf T}=P_3^2,\quad
 P_3\mathbf1=0,\quad \operatorname{tr}P_3=3,
 \qquad P_3P_{11}=P_{11}P_3=P_3.
\]
For the record \eqref{exc:k}, the vector
\begin{equation}
 u_K=P_3P_{11}K
 \label{eq:k-vector-direccional}
\end{equation}
is nonzero. Its exact norm is
\begin{equation}
 \lVert u_K\rVert^2
 =\frac{6638585+2275584\sqrt5}{20}>0.
 \label{eq:k-norma-direccion}
\end{equation}
\end{lemma}
\begin{proof}
The character sum in \eqref{eq:k-proyectores-incidencia} is the isotypic
projector. The real character and the substitution \(a\mapsto a^{-1}\) give
the symmetry. The multiplicity of the sector in the action on \(A_5/C_5\)
is the dimension of its vectors fixed by \(C_5\),
\[
 \frac15\left(3+2\frac{1+\sqrt5}{2}
                 +2\frac{1-\sqrt5}{2}\right)=1.
\]
Therefore its rank is three. Its orthogonality to the trivial character gives
\(P_3\mathbf1=0\) and the two compositions with \(P_{11}\). These
identities may also be verified by multiplying the matrices in the finite
sum.

Finally, \(\sum K_p=6263\), and substituting the twelve components into the
same sum gives
\[
 \lVert u_K\rVert^2=K^{\mathsf T}P_3K
 =\frac1{20}\sum_{a\in A_5}
                \chi_3(a^{-1})K^{\mathsf T}\rho(a)K
 =\frac{6638585+2275584\sqrt5}{20}.
\]
The evaluation of the sixty terms can be checked within the argument itself.
For \(a\in A_5\), set \(q(a)=K^{\mathsf T}\rho(a)K\). Enumerating the even
permutations by cycle type produces the following multisets; in the last
three classes, every printed value has multiplicity two:
\begin{center}
\footnotesize
\setlength{\tabcolsep}{5pt}
\begin{tabular}{@{}c|c|l@{}}
\textbf{class}&\(\lvert C\rvert\)&
\(\{q(a):a\in C\}\), with multiplicities\\ \hline
\(1\)&1&\(4251257\)\\
\(2^2\)&15&
\(2539906,\ 2552978,\ 2664658,\ 2740810\)\\
&&\(2784738,\ 2892172,\ 2911830,\ 3133488\)\\
&&\(3239012,\ 3263782,\ 3404106,\ 3501032\)\\
&&\(3554410,\ 3675796,\ 3809354\)\\
\(3\)&20&
\(2835944,\ 2857732,\ 2943734,\ 3003237\)\\
&&\(3232885,\ 3317239,\ 3354996\)\\
&&\(3410055,\ 3455630,\ 3638920\)\quad
   (each twice)\\
\(5_a\)&12&
\(3396405,\ 3515761,\ 3555976\)\\
&&\(3587516,\ 3592629,\ 3765948\)\quad
   (each twice)\\
\(5_b\)&12&
\(2829137,\ 2917179,\ 3001518\)\\
&&\(3202919,\ 3232872,\ 3955026\)\quad
   (each twice).
\end{tabular}
\end{center}
The table is obtained without approximation: for each of the sixty image
lists \(a\), one locates \(ar_pC_5=r_qC_5\) and sums \(K_pK_q\). Their
class sums are, respectively,
\[
 S_1=4251257,\quad S_{2^2}=46668072,\quad S_3=64100744,\quad
 S_{5_a}=42828470,\quad S_{5_b}=38277302.
\]
Since the character vanishes on the class of order three,
\begin{align*}
 K^{\mathsf T}P_3K
 &=\frac1{20}\left(3S_1-S_{2^2}
       +\frac{1+\sqrt5}{2}S_{5_a}
       +\frac{1-\sqrt5}{2}S_{5_b}\right)\\
 &=\frac{1327717}{4}+\frac{568896}{5}\sqrt5
  =\frac{6638585+2275584\sqrt5}{20}.
\end{align*}
The two positive coefficients prove nonvanishing, and the table permits the
finite evaluation to be repeated directly.
\end{proof}

\begin{definition}[Dodecaphase direction and transverse complement]
The direction selected by \(K\) in the declared chart is
\(\ell_K=\RR u_K\). Its transverse complement within
\(V_{11}=\mathbf1^\perp\) is
\[
 V_{10}(K)=V_{11}\cap u_K^\perp.
\]
\end{definition}

\begin{proposition}[Orthogonal reduction determined by \(K\)]
\label{prop:k-reduccion-diez}
The operator
\begin{equation}
 P_{10}(K)=P_{11}
       -\frac{u_Ku_K^{\mathsf T}}{\lVert u_K\rVert^2}
 \label{eq:k-proyector-diez}
\end{equation}
is the rank-ten orthogonal projector onto \(V_{10}(K)\). One has
\[
 V_{11}=\ell_K\mathbin{\oplus^{\perp}}V_{10}(K),
 \qquad P_{10}P_{11}=P_{10},\quad P_{10}u_K=0.
\]
\end{proposition}
\begin{proof}
Let \(E_K=u_Ku_K^{\mathsf T}/\lVert u_K\rVert^2\). The proved
nonvanishing guarantees that it is defined. It is symmetric,
\(E_K^2=E_K\), has rank one, and, since \(u_K\in V_{11}\),
\(P_{11}E_K=E_KP_{11}=E_K\). Hence
\[
 (P_{11}-E_K)^2=P_{11}-2E_K+E_K=P_{11}-E_K.
\]
Its image is \(V_{11}\cap u_K^\perp\), and its trace is \(11-1=10\). The
decomposition and the remaining identities follow by projection onto the two
orthogonal summands.
\end{proof}

This result specifies a role of \(K\) not exhausted by its designation as a
record: in addition to reversibly representing oriented data and selecting
the tetrad \(B_K\), it determines a linear polarization. The direction
depends on the complete record through \(P_3K\), not merely on the four
points of its support. The causal order is therefore record, representation,
direction, and dimensional reduction.

The transformation has a precise covariance law. If the chart is
simultaneously changed by an orthogonal permutation matrix \(S\), with
\(K'=SK\) and \(P'_3=SP_3S^{-1}\), then
\[
 u_{K'}'=Su_K,\qquad P'_{10}(K')=SP_{10}(K)S^{-1}.
\]
On the other hand, replacing \(u_K\) by \(a u_K\), with \(a\ne0\), does not
change \(E_K\). The direction retains a polarization and does not by itself
permit recovery of either the amplitude of \(u_K\) or the complete record.
The reversibility of \(\kappa\leftrightarrow K\) remains in its proper
domain.

The geometric realization established here is a reduction of inner-product
spaces. Its subsequent use as a compact direction requires transporting this
line to the corresponding physical realization. Section
\ref{exc:pantallas-dualidad} presents the algebraic map that keeps this
reduction and the lattice screen coordinated.
